\notepage{73}{Hydrogen in a uniform electric field}
Let the external electric potential be
\[
 \varphi_{\rm ext}(\mathbf r)=-\mathbf E\cdot\mathbf r,
 \qquad-\boldsymbol\nabla\varphi_{\rm ext}=\mathbf E.
\]
If the electron is at $\mathbf r_1$ and the proton at $\mathbf r_2$,
the perturbation is
\begin{align*}
 H'&=-e\varphi_{\rm ext}(\mathbf r_1)+e\varphi_{\rm ext}(\mathbf r_2)\\
 &=e\bigl[\varphi_{\rm ext}(\mathbf r_2)-\varphi_{\rm ext}(\mathbf r_1)\bigr]\\
 &=e(\mathbf r_1-\mathbf r_2)\cdot\mathbf E
  =e\mathbf r\cdot\mathbf E,
 \qquad\mathbf r\equiv\mathbf r_1-\mathbf r_2.
\end{align*}
In a state of definite orbital parity,
\[
 \mel{n\ell m}{e\mathbf r\cdot\mathbf E}{n\ell m}
 =e\mathbf E\cdot\mel{n\ell m}{\mathbf r}{n\ell m}=0.
\]
Indeed, $\mathbf r$ is an odd-parity operator, so $1+\ell+\ell$ is odd.
It is also a rank-one tensor, with components proportional to $T_1^{m'}$.
\editorial{A zero diagonal matrix element does not rule out a linear
Stark effect in an excited, degenerate hydrogen level. For $n\ge2$, the
perturbation must first be diagonalized within the degenerate subspace.}
\smallheading{Ground-state second-order shift}
For $\mathbf E=E_0\hat{\mathbf e}_z$, $H'=eE_0Z$, and
\[
 E_{100}^2=\sum_{n\ell m\ne100}
 \frac{e^2E_0^2\abs{\mel{n\ell m}{Z}{100}}^2}
      {E_{100}^0-E_{n\ell m}^0}.
\]
Since $Z\propto Y_1^0$ is the $M=0$ component of a vector,
\begin{itemize}
 \item parity requires $\ell$ to be odd;
 \item the triangle rule gives $|0-1|\le\ell\le0+1$, hence $\ell=1$;
 \item the magnetic selection rule gives $m=0+0=0$.
\end{itemize}
For the bound-state terms, this reduces to
\[
 \sum_{n=2}^{\infty}
 \frac{e^2E_0^2\abs{\mel{n10}{Z}{100}}^2}
      {-\mathrm{Ry}\,(1-1/n^2)}.
\]
\editorial{The bound-state sum displayed in the source is only part of
the complete spectral sum: continuum states with the same angular
selection rules also contribute. No continuum evaluation is supplied
on this page.}
